% =====================================================================
%  PAPER I -- merged from proposals P1 (certified PNT substitution) and
%  P2 (the security tax). Target: IEEE S&P, main track.
%
%  Standalone. Set as Overleaf's main document; pdfLaTeX + BibTeX.
%    pdflatex PaperI_certified_budget && bibtex PaperI_certified_budget \
%      && pdflatex PaperI_certified_budget && pdflatex PaperI_certified_budget
% =====================================================================
\documentclass[11pt,a4paper]{article}
% =====================================================================
%  Shared preamble for the three research proposals.
%
%  Each proposal \input{}s this file and is otherwise standalone, so any
%  one of them can be set as Overleaf's main document and compiled on its
%  own. Nothing here is specific to a single proposal.
% =====================================================================
\usepackage[T1]{fontenc}
\usepackage[utf8]{inputenc}
% Times, matching the parent manuscript. mathptmx is in every TeX Live and on
% Overleaf; lmodern is deliberately not used, since it is absent from minimal
% distributions and this preamble should compile anywhere.
\usepackage{mathptmx}
\usepackage[a4paper,margin=25mm]{geometry}
\usepackage{amsmath,amssymb,amsthm}
\usepackage{booktabs}
\usepackage{tabularx}
\usepackage{xcolor}
\usepackage{titlesec}
\usepackage{enumitem}
\usepackage{microtype}
\usepackage[hidelinks,breaklinks]{hyperref}

% ---- shared style tokens, matching the parent manuscript ------------
\definecolor{netcol}{RGB}{31,78,121}
\definecolor{autocol}{RGB}{18,120,90}
\definecolor{bridgecol}{RGB}{176,96,0}
\definecolor{accent}{RGB}{140,20,20}

\newcommand{\layernet}{\textcolor{netcol}{\textbf{network}}}
\newcommand{\layerauto}{\textcolor{autocol}{\textbf{autonomy}}}
\newcommand{\dataset}[1]{\textsc{#1}}
\newcommand{\code}[1]{\texttt{#1}}
\newcommand{\Real}{\mathbb{R}}
\newcommand{\Norm}[1]{\left\lVert #1\right\rVert}
\newcommand{\MCR}{\mathrm{MCR}}

\newtheorem{claim}{Claim}
\newtheorem{proposition}{Proposition}
\newtheorem{definition}{Definition}
\theoremstyle{remark}
\newtheorem{remark}{Remark}

% ---- section styling -------------------------------------------------
\titleformat{\section}{\large\bfseries\color{netcol}}{\thesection}{0.7em}{}
\titleformat{\subsection}{\normalsize\bfseries}{\thesubsection}{0.6em}{}
\titlespacing*{\section}{0pt}{1.4em}{0.6em}

\setlist[itemize]{leftmargin=1.35em,topsep=0.3em,itemsep=0.25em}
\setlist[enumerate]{leftmargin=1.6em,topsep=0.3em,itemsep=0.25em}

% ---- a boxed "what this rests on" callout ---------------------------
\newenvironment{honestly}
  {\par\medskip\noindent\begin{tabular}{@{}p{2.5mm}@{\hspace{3mm}}p{0.9\linewidth}@{}}
   \textcolor{accent}{\rule{2.5mm}{\baselineskip}} &
   \small\itshape}
  {\end{tabular}\par\medskip}

% ---- the proposal header --------------------------------------------
\newcommand{\proposalheader}[3]{%
  \begin{center}
  {\Large\bfseries #1\par}\medskip
  {\large\itshape #2\par}\medskip
  {\small #3\par}
  \end{center}
  \medskip\hrule\medskip}

\usepackage{array}
\newcolumntype{Y}{>{\raggedright\arraybackslash}X}
\newcolumntype{P}[1]{>{\raggedright\arraybackslash}p{#1}}
\newcommand{\Dauth}{\Delta_{\mathrm{auth}}}
\newcommand{\geff}{\gamma_{\mathrm{eff}}}
\newcommand{\gu}{\gamma_{\mathrm{u}}}
\newcommand{\ga}{\gamma_{\mathrm{a}}}
\newcommand{\tfresh}{\tau_{\mathrm{fresh}}}
\newcommand{\eclk}{\varepsilon_{\mathrm{clk}}}
\newcommand{\prov}[1]{\textsc{\small[#1]}}

\begin{document}

\proposalheader
  {Paper I --- The Certified Budget}
  {What Authenticity Costs in Metres: A Certified Budget for Cryptography and
   Sensing in UAV Navigation}
  {Integrated proposal, merging P1 and P2 $\cdot$ RobustUAVs.ai $\cdot$
   Roger Nick Anaedevha $\cdot$ 2026-10-01\\
   Target: IEEE Symposium on Security and Privacy, main track $\cdot$
   follow-on to the parent manuscript~\cite{parentpaper}}

\begin{abstract}
\noindent
A robustness certificate for UAV navigation bounds the position error an
adversary can induce as the product of a delay it can hide, a rate at which that
delay becomes position error, and a dynamics factor. We observe that both of the
first two factors are security-engineering choices, and that a single operation
--- authentication --- moves them in opposite directions. Authenticating a
navigation source lets the estimator trust it, which lowers the error rate
$\gamma$; authenticating a message costs signing, verification and transmission
time, which raises the delay $\Delta$. The certificate is therefore a budget in
which cryptography and sensing trade, and it decides the net sign. We derive the
crossover condition under which authentication increases the certified mission
floor, show that tightening an intrusion detector makes authentication
\emph{more affordable but less beneficial}, and predict that the binding cost of
post-quantum signatures on a flight-controller bus is occupancy, a cliff, rather
than per-message latency, a slope. We propose to measure $\gamma$ for three
navigation sources and the authentication cost of five schemes on
flight-controller-class hardware, to map the crossover across transports, and to
show that an adversary who floods a bus with invalid signatures turns the defence
into the delay attack the certificate already bounds.
\end{abstract}

% =====================================================================
\section{The thesis}
% =====================================================================
The parent manuscript~\cite{parentpaper} certifies a mission-completion floor by
bounding the worst-case position error
\begin{equation}
\label{eq:rho}
\rho(\theta) \;=\; \Delta(\theta)\cdot\gamma\cdot e^{LT_c}\;\le\; m ,
\end{equation}
where $\Delta(\theta)\le H\min(\theta,s)$ is the delay an adversary can hide
below a detector set at $\theta$ across $H$ hops, $\gamma$ is metres of error per
second of degraded input, $L$ is the local Lipschitz constant of the navigation
dynamics, $T_c$ the certification horizon, and $m$ the corridor margin. Two of
the four inputs are fixed by physics ($L$) and regulation ($m$). The other two
are not.

\begin{itemize}
\item $\gamma$ depends on \emph{which navigation source the estimator trusts}.
  P1 observed that it is a property of the source, not of the attack, and that
  an authenticated source can be trusted under conditions where an
  unauthenticated one cannot.
\item $\Delta$ is indifferent to \emph{why} an input is stale. P2 observed that
  authentication itself makes inputs stale, by exactly the quantity the theorem
  already bounds.
\end{itemize}

\noindent
Read together: \textbf{authentication is both a credit and a debit against one
certified budget.} It buys a lower $\gamma$ and pays for it in $\Delta$. Neither
proposal alone could ask which effect wins. This paper asks it, answers it with
a closed-form condition, and measures where real transports and schemes fall.

% =====================================================================
\section{Threat model: what authentication actually buys}
% =====================================================================
We inherit the parent's adversary --- a compromised mesh relay that forwards
every message late, and a GNSS jammer or spoofer --- and must be precise about
what authentication does against each, because the obvious answer is wrong.

\paragraph{Against forgery, authentication works.} A spoofed navigation fix
that fails verification is rejected rather than fused. This is P1's mechanism,
and it is what lowers $\gamma$: the estimator stops absorbing the spoofer's
position offset.

\paragraph{Against delay, authentication alone does nothing.} A validly signed
message forwarded four seconds late is still validly signed. The parent's
time-delay adversary is untouched by signatures. Authentication bounds delay
only when the signature covers a timestamp and the receiver rejects messages
older than a freshness window $\tfresh$. MAVLink~2 signing already carries such a
timestamp~\cite{trout_mavlink}. Freshness then caps hidden delay at
$\tfresh+\eclk$, where $\eclk$ is the receiver's clock error.

\paragraph{The clock couples back to the attack.} On a UAV the reference clock
is very often disciplined by GNSS. Under jamming it free-runs, and $\eclk$ grows
with holdover duration. So the freshness bound weakens precisely when the
navigation attack is active. This coupling is, to our knowledge, unmodelled in
both the drone-authentication and the certified-robustness literature, and a
security reviewer will raise it first. The paper treats it as a first-class
term:
\begin{equation}
\label{eq:delta-eff}
\Delta_{\mathrm{eff}}(\theta,t) \;=\;
\min\!\big(\,H\min(\theta,s),\;\; \tfresh+\eclk(t)\,\big) \;+\; \Dauth .
\end{equation}

\paragraph{Symmetric MACs do not stop a compromised relay.} MAVLink~2's
HMAC-SHA256 uses a shared key. A compromised relay holds that key and can sign
anything. Against the parent's adversary, therefore, only asymmetric signatures
or delayed-disclosure schemes such as TESLA~\cite{perrig2000tesla} provide
source authentication --- and those are exactly the expensive ones. This is why
the cost side of the budget matters and is not a footnote.

\paragraph{New adversary: induced verification.} An adversary on the bus sends
frames carrying invalid signatures. The receiver must spend verification time
before it learns they are invalid, so the attacker converts the defence's own
cost into delay (\S\ref{sec:attack}).

% =====================================================================
\section{The certified budget}
\label{sec:theory}
% =====================================================================
Let $S$ be the set of available navigation sources with authentication states
$a$, and $\kappa$ an authentication scheme operating at message rate $r$ on a
transport with per-frame payload $b$ and frame duration $t_f$. The merged
budget is
\begin{equation}
\label{eq:budget}
\rho(\theta,S,a,\kappa) \;=\;
\Delta_{\mathrm{eff}}(\theta)\;\cdot\;\geff(S,a)\;\cdot\;e^{LT_c}\;\le\;m .
\end{equation}
Write $E=e^{LT_c}$, and in the delay regime, where re-anchoring holds by
construction and the freshness cap is not binding, $\Delta_{\mathrm{eff}} =
\Delta(\theta)+\Dauth$.

\begin{proposition}[Crossover]
\label{pr:crossover}
Let authenticating the navigation path change the effective rate from $\gu$ to
$\ga<\gu$ at cost $\Dauth$. Authentication strictly reduces the certified
error $\rho$ if and only if
\begin{equation}
\label{eq:crossover}
\Dauth \;<\; \Delta(\theta)\left(\frac{\gu}{\ga}-1\right).
\end{equation}
\end{proposition}
\noindent\emph{Proof.} $(\Delta+\Dauth)\,\ga E < \Delta\,\gu E$, divide by
$\ga E>0$ and rearrange. \hfill$\square$

\smallskip
\noindent
The condition is elementary; what is not elementary is that its two sides are
measured on entirely different instruments --- a logic analyser on a CAN bus and
a flight log of a jammed receiver --- and that nobody has put them in one
inequality.

\begin{proposition}[More affordable, less beneficial]
\label{pr:afford}
Let $B(\theta)=m/(\ga E)-\Delta(\theta)$ be the absolute authentication budget
and $C(\theta)=\Delta(\theta)(\gu/\ga-1)$ the crossover threshold of
Eq.~\eqref{eq:crossover}. For $\theta<s$, $B$ is decreasing in $\theta$ and $C$
is increasing in $\theta$. Tightening the detector therefore leaves
\emph{more} room to afford authentication and \emph{less} reason to want it.
\end{proposition}
\noindent\emph{Proof.} $\Delta(\theta)=H\theta$ on $\theta<s$; $B$ has slope
$-H$ and $C$ has slope $H(\gu/\ga-1)>0$. \hfill$\square$

\smallskip
\noindent
This refines P2's first proposition, which stated only the affordability half.
The two halves answer different questions --- \emph{can I pay for it?} and
\emph{is it worth paying for?} --- and an engineer who conflates them will
over-invest in cryptography on a platform whose detector is already tight.

\begin{proposition}[The cost is a cliff, not a slope]
\label{pr:cliff}
$\Dauth$ decomposes into compute, transmission and protocol terms,
\[
\Dauth \;=\; \underbrace{t_{\mathrm{sign}}+t_{\mathrm{verify}}}_{\text{compute}}
\;+\; \underbrace{\lceil\sigma/b\rceil\,t_f\cdot q(U)}_{\text{transmission}}
\;+\; \underbrace{t_{\mathrm{proto}}}_{\text{protocol}} ,
\]
where $\sigma$ is the tag size and $q(U)$ a queueing factor in bus occupancy
$U=U_0+r\lceil\sigma/b\rceil t_f$. Since $q(U)\to\infty$ as $U\to1$, every scheme
has a maximum admissible message rate
$r^{*}=(1-U_0)/(\lceil\sigma/b\rceil t_f)$, above which $\Dauth$ is unbounded and
the certificate is vacuous regardless of $m$. Ordering schemes by $\sigma$
captures only the transmission term: schemes that minimise it, such as TESLA,
can maximise $t_{\mathrm{proto}}$ through their key-disclosure delay.
\end{proposition}

\smallskip
\noindent
This is P2's third proposition made precise, and it is the one we expect to
carry the empirical weight (\S\ref{sec:already}).

\begin{proposition}[Two regimes]
\label{pr:regimes}
Under the time-delay adversary, re-anchoring holds by construction and the
crossover of Proposition~\ref{pr:crossover} decides. Under a
\emph{persistent-denial} adversary that withholds every unauthenticated source
for longer than $T_c$, the unauthenticated certificate is vacuous, and
authentication of at least one source is \emph{necessary} for any certified
floor.
\end{proposition}

\smallskip
\noindent
This is P1's second claim, stated as a regime boundary rather than a universal.

\paragraph{The selection rule.} Given $m$, $\theta$, a transport, and a cost per
source and scheme, Eq.~\eqref{eq:budget} induces a feasibility set over
(suite, scheme) pairs; the rule returns its cost-minimal element. This is P1's
third claim, extended to choose the scheme as well as the suite. It is verified
against brute-force enumeration; the optimality argument goes to the appendix.

% =====================================================================
\section{What the parent's constants already say}
\label{sec:already}
% =====================================================================
Two things can be computed today, before any new measurement, from constants
the parent manuscript already reports: $H=2$, $\theta=0.25$\,s, $L=1.181$,
$T_c=1$\,s, and the campaign supremum $\gamma=1.625$\,m/s. They reproduce the
parent's $\gamma_{\mathrm{req}}=6.14$\,m/s at $m=10$\,m, which is the check that
the arithmetic below is wired to the same numbers.

\begin{center}\small
\begin{tabular}{@{}rccc@{}}
\toprule
$m$ (m) & window edge $\theta^{*}$ (s) & certified at $\theta=0.25$\,s? &
authentication budget $B$ (s) \\
\midrule
2  & 0.189 & no  & $-0.122$ \\
5  & 0.472 & yes & $\phantom{-}0.445$ \\
10 & 0.945 & yes & $\phantom{-}1.389$ \\
20 & 1.889 & yes & $\phantom{-}3.278$ \\
\bottomrule
\end{tabular}

\smallskip
\prov{derived} from measured parent constants; $B$ uses $\ga=\gu=1.625$\,m/s,
i.e.\ assumes authentication buys no $\gamma$ improvement.
\end{center}

\noindent
Two readings follow. First, at $m=2$\,m nothing certifies \emph{even at zero
authentication cost}: P2's second proposition is already partly a consequence
of the parent's results, and the new content is naming that margin per transport.
Second, at $m=5$\,m the budget is $0.445$\,s. Per-message signing latency on a
microcontroller is measured in milliseconds, so \textbf{per-message latency is
very unlikely to be the binding constraint.} Something else must be.

\paragraph{Order-of-magnitude motivation for the cliff.}
The table below is a back-of-envelope calculation, not a result, and is labelled
as such. Tag sizes are from the standards~\cite{nist_pqc}; the frame model
assumes classic CAN at 1\,Mbit/s, seven payload bytes per DroneCAN multi-frame
transfer frame, and roughly 130 bits per frame including stuffing.

\begin{center}\small
\begin{tabular}{@{}lrrrr@{}}
\toprule
Scheme & tag (B) & frames & ms / message & bus saturates at (msg/s) \\
\midrule
HMAC-SHA256 (MAVLink~2 signature field) & 13   & 2    & 0.26  & $\sim$3800 \\
Ed25519                                 & 64   & 10   & 1.3   & $\sim$770 \\
ML-DSA-44 (FIPS 204)                    & 2420 & 346  & 45    & $\sim$22 \\
SLH-DSA-SHA2-128s (FIPS 205)            & 7856 & 1123 & 146   & $\sim$7 \\
\bottomrule
\end{tabular}

\smallskip
\prov{estimate} --- to be replaced by measurement (\S\ref{sec:measure}).
\end{center}

\noindent
If these estimates are even roughly right, an ML-DSA-44 signature costs about
$45$\,ms per message --- comfortably inside a $445$\,ms budget --- yet saturates a
classic CAN bus at about 22 signed messages per second, below the update rate of
many attitude and inertial streams. That is Proposition~\ref{pr:cliff} in
numbers: the scheme is affordable per message and inadmissible per second.

\begin{honestly}
These are the most quotable numbers in the proposal and the least established.
They rest on an assumed frame duration and an assumed per-frame payload, and
real buses carry a base load $U_0$ that only lowers the saturation rates
further. They motivate the measurement; they must never appear in the paper
except as the prediction the measurement tested.
\end{honestly}

% =====================================================================
\section{Research questions}
% =====================================================================
\begin{description}[leftmargin=2.6em,style=nextline]
\item[RQ1] \textbf{How much does $\gamma$ depend on the navigation source?}
  Measured for GNSS (done: $1.625$\,m/s), inertial-only, and visual--inertial.
\item[RQ2] \textbf{What does authentication cost on real flight hardware?}
  Compute, transmission and protocol terms of $\Dauth$, for five schemes, two
  targets, three transports.
\item[RQ3] \textbf{Where is the crossover?} For each (transport, scheme,
  source) triple, the sign of Eq.~\eqref{eq:crossover} and the margin at which it
  flips.
\item[RQ4] \textbf{Does the freshness bound survive GNSS denial?} The growth of
  $\eclk$ under holdover, and the denial duration after which
  Eq.~\eqref{eq:delta-eff} reverts to the detector bound.
\item[RQ5] \textbf{How much delay can an attacker induce through
  verification?} And how much of it does a pre-filter recover?
\end{description}

% =====================================================================
\section{Measurement plan}
\label{sec:measure}
% =====================================================================
\subsection{$\gamma$ per navigation source (RQ1)}
\begin{enumerate}
\item \textbf{Inertial-only --- first, and this week.} From the Whelan flight
  logs already on disk~\cite{whelan2020uav}: freeze the GNSS fix at onset and
  integrate the inertial measurements. This replaces an error-budget
  \emph{assertion} in the parent's cruise discussion with a \emph{measurement},
  so it strengthens the parent paper as well as this one.
\item \textbf{Visual--inertial.} Public VIO datasets with UAV motion profiles,
  using the parent's baseline-corrected estimator $\gamma_\delta$ rather than the
  cumulative one, which the parent showed is contaminated by detector
  threshold.
\item \textbf{LEO signals of opportunity --- appendix only.} No receiver logs
  are available. A \code{physics\_model} value derived from published
  RMSE~\cite{kassas2021leodoppler,kassas2022jammer} is reported, tagged, and kept
  out of every headline number.
\end{enumerate}

\subsection{$\Dauth$ per scheme, target and transport (RQ2)}
\begin{itemize}
\item \textbf{Schemes:} HMAC-SHA256 (MAVLink~2 baseline, shared key), Ed25519,
  ML-DSA-44, SLH-DSA-SHA2-128s~\cite{nist_pqc}, and TESLA-style delayed-key
  disclosure~\cite{perrig2000tesla}, whose $t_{\mathrm{proto}}$ is an explicit
  tunable parameter and so gives the cleanest test of
  Proposition~\ref{pr:cliff}.
\item \textbf{Targets:} an STM32H7-class microcontroller (as on Pixhawk~6) and a
  companion-computer-class ARM board. Cycle-accurate timing on the first,
  wall-clock under load on the second.
\item \textbf{Transports:} classic CAN (DroneCAN), CAN-FD, and MAVLink~2 over a
  serial telemetry link.
\item \textbf{Validation:} the fragmentation model is checked against inter-frame
  timing of the benign portion of the \dataset{HCRL UAVCAN}
  capture~\cite{hcrl_uavcan}, which also supplies a realistic base load $U_0$.
\item \textbf{Amortisation as a design axis:} per-message signing, batch signing
  over $k$ messages, and MAC-plus-periodic-signature hybrids each trade
  bandwidth for latency, which is precisely the budget's trade.
\end{itemize}

\subsection{Clock holdover (RQ4)}
Determine first whether the released flight logs carry both a GNSS-disciplined
time and a free-running system time. If they do, $\eclk(t)$ under the recorded
jamming intervals is directly measurable. If they do not, use oscillator
holdover specifications for the flight-controller crystal, tagged
\code{physics\_model}.

\begin{honestly}
We do not yet know whether the Whelan logs carry both clocks. This must be
checked before RQ4 is promised, and if they do not, RQ4 becomes a modelled
analysis and the paper says so.
\end{honestly}

% =====================================================================
\section{The induced-verification attack (RQ5)}
\label{sec:attack}
% =====================================================================
An adversary with bus access injects frames carrying syntactically valid but
cryptographically invalid signatures at rate $r_{\mathrm{adv}}$. Each costs the
receiver $t_{\mathrm{verify}}$ before rejection, and on a multi-frame scheme it
also costs reassembly and bus occupancy. The induced delay on legitimate
navigation traffic enters Eq.~\eqref{eq:delta-eff} exactly as adversarial
staleness does, so the composition theorem bounds the mission-level effect.

\paragraph{Positioning against prior art.} Signature-flooding denial of service
is not new: avoiding per-packet signature verification was a design motivation
for TESLA~\cite{perrig2000tesla}. The contribution is not the mechanism. It is
(i) the \emph{mission-level bound} on what the attack achieves, through the
certificate, and (ii) the observation that post-quantum migration makes it
markedly worse on constrained buses, because lattice verification is costlier
than the classical verification it replaces and multi-frame reassembly must
complete before verification can begin.

\paragraph{Mitigation in the same section.} Cheap pre-filters (a truncated MAC
checked before the full signature), per-source rate limits, and verification
budgets, each evaluated by how much of the induced $\Delta$ it recovers. As in
the parent's evasion analysis, the output is a defensive configuration
procedure, and the attack is presented alongside it rather than ahead of it.

% =====================================================================
\section{Contributions, as they would appear in the paper}
% =====================================================================
\begin{itemize}
\item \textbf{A certified budget} in which sensing and cryptography trade
  through a single robustness certificate, with a closed-form crossover for when
  authentication helps (Prop.~\ref{pr:crossover}), the affordability--benefit
  inversion (Prop.~\ref{pr:afford}), and a two-regime boundary
  (Prop.~\ref{pr:regimes}).
\item \textbf{The freshness--clock coupling}: an explicit model of how GNSS
  denial weakens the freshness check that is authentication's only defence
  against delay.
\item \textbf{Measurements} of $\gamma$ for three navigation sources and of
  $\Dauth$ for five schemes on flight-controller-class hardware across three
  transports, every number provenance-tagged.
\item \textbf{The cliff}: evidence that bus occupancy, not per-message latency,
  determines which post-quantum schemes a UAV can adopt
  (Prop.~\ref{pr:cliff}).
\item \textbf{An attack with its bound and its mitigation}: induced verification,
  bounded at the mission level.
\item \textbf{A selection rule} returning the cheapest (sensor suite, scheme)
  pair that keeps a given detector setting certified.
\end{itemize}

% =====================================================================
\section{Relation to prior work}
% =====================================================================
\begin{center}\small
\begin{tabularx}{\linewidth}{@{}P{0.28\linewidth}YY@{}}
\toprule
\textbf{Line of work} & \textbf{What it reports} & \textbf{What this paper adds} \\
\midrule
Post-quantum authentication for drones~\cite{pq_drone_networks,utm_survey2026}
& signature size, cycles, energy & what those milliseconds do to the aircraft \\
Broadcast authentication (TESLA)~\cite{perrig2000tesla} & cheap source
authentication at the price of disclosure delay & that delay priced in metres,
against a certified margin \\
Alternative PNT~\cite{kassas2021leodoppler,kassas2022jammer} & positioning
accuracy under denial & accuracy translated into a certified window \\
Certified robustness~\cite{cohen2019smoothing,gronwall1919} & stability under
given perturbations & where the perturbation comes from, and what defending
against it costs \\
The parent manuscript~\cite{parentpaper} & a certificate with $\gamma$ and
$\Delta$ as fixed inputs & both as design variables, jointly optimised \\
\bottomrule
\end{tabularx}
\end{center}

% =====================================================================
\section{Relation to the parent manuscript's limitations}
% =====================================================================
\begin{center}\small
\begin{tabularx}{\linewidth}{@{}P{0.40\linewidth}P{0.10\linewidth}Y@{}}
\toprule
\textbf{Limitation} & \textbf{Closed?} & \textbf{How} \\
\midrule
The certificate assumes re-anchoring & \textbf{fully} & Prop.~\ref{pr:regimes}
makes it a regime boundary with a named satisfying mechanism. \\
$\gamma$ from one platform in one regime & \textbf{partly} & Adds
sources; does \emph{not} deliver a cruise campaign, which stays data-gated. \\
One attack surface & \textbf{partly} & Adds the bus and the C2 link on the
defence side. \\
Evidence-class imbalance & no & Unchanged. \\
\bottomrule
\end{tabularx}
\end{center}

% =====================================================================
\section{Risks and fallbacks}
% =====================================================================
\begin{itemize}
\item \textbf{The crossover could be one-sided.} If every measured $\Dauth$ is
  far below $C(\theta)$, authentication always wins and Proposition~\ref{pr:crossover}
  has no interesting boundary. The cliff makes this unlikely on classic CAN, but
  if it happens the paper reports it plainly: ``authentication always pays,
  except where occupancy forbids it'' is still a result, and the occupancy
  boundary becomes the headline.
\item \textbf{Hardware access.} A development board is inexpensive; a full flight
  controller is needed only for bus validation. Fallback: a pinned ARM target,
  tagged as such. The crossover survives because it depends on the ratio
  $\Dauth/\Delta(\theta)$, which a pinned target still bounds.
\item \textbf{Only two $\gamma$ sources measured.} GNSS and inertial-only already
  establish that $\gamma$ is source-dependent. VIO strengthens; it is not
  load-bearing.
\item \textbf{The freshness coupling might dominate.} If $\eclk$ grows fast
  enough under holdover, freshness provides no protection within the
  certification horizon and the paper's message becomes ``authenticate the
  source, not the timing''. That is a valid and useful outcome.
\item \textbf{Scheme selection dates.} Pin to FIPS 203/204/205 and publish the
  model, so a new scheme can be placed without re-running the campaign.
\end{itemize}

% =====================================================================
\section{Page plan and timeline}
% =====================================================================
\begin{center}\small
\begin{tabular}{@{}lr@{\qquad}lr@{}}
\toprule
\textbf{Section} & \textbf{pp} & \textbf{Section} & \textbf{pp} \\
\midrule
Introduction                    & 1.25 & $\Dauth$ results and the cliff & 1.5 \\
Background and threat model     & 1.25 & Crossover map and selection    & 1.25 \\
The certified budget            & 2.0  & Induced verification           & 1.25 \\
Measurement methodology         & 1.5  & Discussion and limitations     & 0.75 \\
$\gamma$ results                & 1.25 & Related work, conclusion       & 1.0 \\
\midrule
\multicolumn{3}{@{}l}{\textbf{Body total}} & \textbf{13.0} \\
\bottomrule
\end{tabular}
\end{center}

\noindent
About five months. Weeks 1--2: inertial $\gamma$ and the budget calculator
(both need nothing new). Weeks 3--10: hardware benchmarks, VIO, fragmentation
validation. Weeks 8--12: the attack and its mitigation. Weeks 13--20: writing,
with Paper~II submitted first so that its vocabulary is citable.

\small
\bibliographystyle{plain}
\bibliography{proposals}

\end{document}
